\section{Lexical and file syntax}

\subsection{Logical lines, whitespace, and case}

Commands and flags are case-insensitive and are separated by whitespace.
Token names retain their spelling. CRLF and LF line endings are accepted.
A physical line ending \emph{exactly} in a backslash joins to the next line:

\begin{lstlisting}
PARAMS: A:state B:state \
        Cin:state
CCNOT -I A B \
      -O Carry
\end{lstlisting}

The continuation backslash must be the final character; spaces after it prevent
continuation. Blank logical lines are discarded.

\subsection{Hash comments}

\begin{lstlisting}
# Whole-line or trailing hash comment
H -I Q:0 -O Q:0  # prepare |+>
\end{lstlisting}
A hash starts a comment that continues through the logical line. Comments are
for explanations and do not become instructions. Because continuation joining
happens first, a hash on a continued logical line can hide text that appeared
on a later physical line. Prefer comments after a complete instruction.

\subsection{Block comments}

\begin{lstlisting}
/* A block comment may span logical lines. */
PHASE=pi/2 -I Q:0 -O Q:0
\end{lstlisting}

Text from \code{/*} through \code{*/} is ignored. A block may occupy one line
or continue across several lines. Block comments are useful for temporarily
annotating a region, but they should not be nested. Keep delimiters visually
obvious so an omitted closing marker does not hide the remainder of a file.

\subsection{\code{PARAMS:}: the process input declaration}

\begin{lstlisting}
PARAMS: A:state B:state Theta:float Count:int
\end{lstlisting}

\code{PARAMS:} is recognized only when it is the first non-comment logical
line. Each whitespace-separated entry containing a colon becomes
\code{name:type}. Declaration order is interface order: RUNCHILD inputs and
truth-table input columns map from left to right.

Long declarations can use continuation:
\begin{lstlisting}
PARAMS: A0:state A1:state B0:state B1:state \
        Cin:state Width:int
\end{lstlisting}

The system retains the type text:
\begin{itemize}
  \item \code{state} entries become user-facing circuit inputs and truth-table
  input columns. Each participating state parameter resolves to a qubit wire.
  \item \code{int} and \code{float} entries are retained but excluded from
  user-facing qubit controls and truth-table state inputs. They describe
  numeric process metadata rather than prepared qubits.
  \item other type names are retained as written. Because they are not one of
  the recognized numeric exclusions, they may be exposed like stateful
  parameters; prefer \code{state}, \code{int}, or \code{float} for predictable
  behavior.
\end{itemize}

PARAMS has a colon after the keyword; \code{MAIN-PROCESS} does not. An empty or
misplaced PARAMS line is not a substitute for a valid first-line declaration.
The process's MAIN-PROCESS instruction is described in the composition chapter.

\subsection{Protocol file forms}

The normal authoring format is plain UTF-8 QPU text saved as
\code{ProcessName.qpucir}. A \code{.qpucir.txt} download is also plain source
and is useful where browsers or editors insist on a text extension. Uploading
plain text derives the display name from MAIN-PROCESS.

The loader can also accept a JSON export envelope:
\begin{lstlisting}[language={}]
{
  "format": "qpucir",
  "version": 1,
  "name": "First circuit",
  "source": "MAIN-PROCESS FirstCircuit\nRETURNVALS Q"
}
\end{lstlisting}
The source string is the authoritative language text. Export envelopes may
also contain compiled qubit count, gates, token map, and export timestamp
metadata, but compilation from source determines current behavior.

